% Datasets, evaluation protocol and comparability: text for the paper (#138).
% Citation keys are in artifacts/issue-138/references.bib. Every number is
% taken from artifacts/issue-138/reference_statistics.md, release_checks.md
% and artifacts/issue-137/evaluation_protocol.md; the per-benchmark detail
% behind the last subsection is in benchmark_protocols.md.
% Requires \usepackage{threeparttable} (already needed by Table 6).

\subsection{Datasets and reference alignments}
\label{sec:datasets}

We use five datasets (Table~\ref{tab:datasets}), all English. D1 is the
SNOMED~CT--FMA \emph{Body} equivalence task of Bio-ML~2022~\cite{he2022bioml}.
Of its 7,256 reference mappings, the unsupervised setting holds out 726 for
validation and 6,530 for testing; we evaluate on the 6,530 test mappings and
treat the 726 validation mappings as known pairs (see below). D2 is the OAEI
Anatomy task, mouse anatomy against the anatomy part of the NCI
Thesaurus~\cite{bodenreider2005mice,oaei2024anatomy}. D3 comprises the 21
pairs of the seven OntoFarm conference ontologies used by the OAEI Conference
track~\cite{zamazal2017ontofarm,oaei2024conference}, with the original
reference alignment ra1, including both class and property correspondences
(variant M3). D4 is the memoryalpha--stexpanded test case of the OAEI
Knowledge Graph track~\cite{hertling2020kgtrack}, whose gold standard is
partial; we evaluate its schema (class and property) and instance
correspondences as two separate pairs, D4$_{\mathrm{sc}}$ and
D4$_{\mathrm{in}}$. D5 is the English DBpedia--Wikidata 15K dataset (V2) of
the OpenEA benchmark~\cite{sun2020openea}, fold~1, whose entity labels were
removed by its authors.

All input files are byte-identical to the published releases: D1--D4 to the
OAEI track repository distributed through MELT~\cite{hertling2019melt}, D3's
references additionally to the Conference track's ra1 archive (all 21 contain
the same correspondences and ontology declarations; in one, the MELT copy
moves a wrapper element that the original misplaces), and D5 to the
OpenEA v2.0 archive, whose checksum matches the published one. Every
reference correspondence is an equivalence with confidence 1.0; none was
filtered, normalised or deduplicated. The D4 schema/instance partition, by the
source entity's \texttt{rdf:type}, reproduces the track's published sizes
(13 class, 41 property and 1,725 instance correspondences). Reference
alignments are not always one-to-one: in D1, 1,477 of the 6,530 mappings
belong to a source with several targets, and 563 of the 4,605 D1 test
sources have two or more gold targets.

\begin{table*}[t]
\scriptsize
\setlength{\tabcolsep}{4pt}
\begin{threeparttable}
\caption{Datasets, reference alignments and evaluation populations. Refs:
validation/test reference correspondences used here (D1--D4: the reference
split 20/80 by source entity; D5: OpenEA fold~1). Test src: test sources
with at least one reference. Targets: entities stage~1 ranks for every
source. R@20: stage-1 candidate recall@20 on test references, the ceiling of
every re-ranking result (D3: mean over 21 pairs).}
\label{tab:datasets}
\begin{tabular}{llllrrrr}
\hline
\textbf{Col} & \textbf{Task} & \textbf{Release} & \textbf{Reference} & \textbf{Refs (val/test)} & \textbf{Test src} & \textbf{Targets} & \textbf{R@20} \\
\hline
D1 & SNOMED$\to$FMA, classes & Bio-ML 2022 & unsup.\ test, 6,530 & 1,304 / 5,226 & 4,605 & 64,726 & 0.858 \\
D2 & MA$\to$NCI, classes & OAEI Anatomy & full, 1,516 & 302 / 1,214 & 1,198 & 3,304 & 0.947 \\
D3 & 21 pairs, classes+props & OAEI Conference & ra1-M3, 305 & 54 / 251 & 250 & 74--181 & 0.862 \\
D4$_{\mathrm{sc}}$ & memoryalpha$\to$stexp., schema & OAEI KG v4 & partial, 54 & 10 / 44 & 44 & 1,947 & 0.932 \\
D4$_{\mathrm{in}}$ & memoryalpha$\to$stexp., instances & OAEI KG v4 & partial, 1,725 & 345 / 1,380 & 1,380 & 39,305 & 0.973 \\
D5 & DBpedia$\to$Wikidata, instances & OpenEA v2.0 & fold 1, 15,000 & 1,500 / 10,500 & 10,500 & 14,992 & 0.267 \\
\hline
\end{tabular}
\end{threeparttable}
\end{table*}

\subsection{Evaluation protocol}
\label{sec:protocol}

Every condition re-ranks the same candidate lists. Stage~1 ranks all target
entities of the pair's entity type for each source by character 3-gram cosine
over labels (on D5, over attribute values, since labels are absent) and keeps
the top 20. A condition verbalises the source and its 20 candidates and
re-orders them by embedding cosine. As a reference point, we also score the
unchanged stage-1 order (B0). Because candidates, populations, thresholds and
metric definitions are identical for every condition and for B0, differences
between them isolate the effect of re-ranking.

For D1--D4, each reference alignment is split by source entity into 20\%
validation and 80\% test, with a fixed seed, so that all correspondences of a
source fall into the same split. D5 uses OpenEA's fold~1 as published (3,000
train, 1,500 validation, 10,500 test links). No model is trained on any
reference. Every source is scored in exactly one population: test if it has
a test reference, otherwise validation if it has a validation reference;
sources whose only references are known pairs (D5's training links and D1's
Bio-ML validation mappings) are excluded, and known pairs are removed from
every candidate list, following Bio-ML's rule that reference mappings outside
the evaluated set are excluded from system output~\cite{he2022bioml}. Sources
with no reference are divided between the two populations in proportion to
their reference sources. Under a complete reference (D1, D2, D3, D5) a
prediction for such a source is a false positive. Under D4's partial
reference it is judged only when its target is a gold target, mirroring the
track's rule that a correspondence is wrong only if it contradicts a gold
mapping~\cite{hertling2020kgtrack}.

Each source predicts only its top-ranked candidate, kept if its score reaches
a threshold (the one-mapping-per-source prediction of
BERTMap~\cite{he2022bertmap}; BERTMap itself is not run). The threshold is
chosen per pair on the validation population only, as the F1-maximising
value on a grid from 0.10 to 0.99 in steps of 0.01, and applied unchanged to
the test population.

\subsection{Metrics}
\label{sec:metrics}

We report ranking and matching metrics separately. \emph{Ranking} metrics
need no threshold and are computed within the 20 candidates. MRR is the mean,
over test reference sources, of the reciprocal rank of the source's
best-ranked gold target; a source whose gold targets are all missing from its
list, or that has no list, contributes 0. Recall@$k$ is the fraction of test
reference \emph{correspondences} whose target is in the source's top $k$.
\emph{Matching} metrics are precision, recall and F1 of the thresholded
top-1 predictions on the test population; recall is over all test
references, including those whose source was never ranked (167 on D5).
Because each source predicts one target, a source with several gold targets
can recover only one of them. On D1 this rule alone caps recall at 0.881
(4,605 test sources for 5,226 test references); together with the stage-1
candidate lists, which miss every gold target of 525 test sources, the cap
is 0.781. Each column reports one pair, except D3, which reports
the unweighted mean over its 21 pairs. Statistical tests treat the dataset
columns as observations, since D3's pairs share ontologies and are not
independent.

\subsection{Relation to benchmark protocols}
\label{sec:comparability}

This is a controlled comparison under one re-ranking protocol, not a
reproduction of any benchmark's evaluation, and our numbers are not directly
comparable with published system results. The benchmarks evaluate matching
over the full ontologies or knowledge graphs without reference splits:

\begin{itemize}
\item \textbf{Bio-ML} computes global matching P/R/F1 on all 6,530 test
  mappings and local ranking against 100 sampled negatives plus the gold
  target per mapping~\cite{he2022bioml}.
\item \textbf{OAEI Anatomy} scores all 1,516 correspondences and reports
  recall on non-trivial correspondences~\cite{oaei2024anatomy}.
\item \textbf{OAEI Conference} reports micro-averaged results against the
  revised reference rar2~\cite{oaei2024conference}.
\item \textbf{The KG track} matches schema and instances jointly over the
  whole graphs~\cite{hertling2020kgtrack}.
\item \textbf{OpenEA} trains supervised models on 20\% of the links and ranks
  each test source against all test targets, averaged over five
  folds~\cite{sun2020openea}.
\end{itemize}

We instead evaluate 80\% of each reference after tuning a threshold on the
other 20\%, rank within 20 stage-1 candidates (so that the gold target is
absent for up to 73\% of test references on D5, Table~\ref{tab:datasets}),
use ra1 rather than rar2 on D3 and macro-average its pairs, and train
nothing on D5. We therefore compare conditions with each other and with B0
on identical inputs, and do not rank them against published scores.
